% Generated draft: requires mathematical and source review.
\documentclass[11pt]{article}
\usepackage[margin=25mm]{geometry}
\usepackage{fontspec}
\setmainfont{Latin Modern Roman}
\usepackage{amsmath,amssymb,mathtools,unicode-math}
\setmathfont{Latin Modern Math}
\usepackage{enumitem,needspace,xurl,hyperref}
\setlength{\parindent}{0pt}
\setlength{\parskip}{5pt}
\setlength{\emergencystretch}{3em}
\newcommand{\N}{\mathbb{N}}
\newcommand{\Z}{\mathbb{Z}}
\newcommand{\Q}{\mathbb{Q}}
\newcommand{\R}{\mathbb{R}}
\newcommand{\C}{\mathbb{C}}
\newcommand{\F}{\mathbb{F}}
\newcommand{\A}{\mathbb{A}}
\newcommand{\PP}{\mathbb{P}}
\newcommand{\E}{\mathbb{E}}
\newcommand{\eps}{\varepsilon}
\newcommand{\dd}{\,\mathrm{d}}
\begin{document}
\section*{UnsolvedMath ID: 2517}
\textbf{Source problem code:} catalogue entry
\subsection{Definitions and mathematical statement}
For an integer $m\geq1$ and $0\leq r<m$, write $r(m)=\{r+tm:t\in\Z\}$. If $r_1(m_1)$ and $r_2(m_2)$ are disjoint, define the class transposition $\tau_{r_1(m_1),r_2(m_2)}\in\operatorname{Sym}(\Z)$ by
\[
r_1+tm_1\longleftrightarrow r_2+tm_2\qquad(t\in\Z),
\]
fixing every other integer. Disjointness makes this an involution. Let $\operatorname{CT}(\Z)$ be the subgroup generated by all such transpositions.

For $k\geq2$, let $\operatorname{CT}_k$ be generated by the horizontal class transpositions $\tau_{r_1(k),r_2(k)}$ with $0\leq r_1\ne r_2<k$. These permute the $k$ residue labels within each block and give $\operatorname{CT}_k\cong S_k$. Distinctly, set
\[
\operatorname{CT}(k)=\langle\operatorname{CT}_2,\operatorname{CT}_3,\ldots,\operatorname{CT}_k\rangle,
\qquad N=\operatorname{lcm}(2,3,\ldots,k).
\]
The question is whether, for every integer $k>3$,
\[
\operatorname{CT}(k)\cong S_N.
\]
The supplied assessment identifies the intended action more specifically as all permutations of the $N$ residue labels in the canonical blocks $qN+\{0,\ldots,N-1\}$, acting the same way in every block. Thus the subgroup on these labels is asserted to be the full symmetric group, not merely some subgroup or a permutation of the infinite set with unspecified action. The subscript notation $\operatorname{CT}_k$ and parenthesized $\operatorname{CT}(k)$ are not interchanged.

The source is Kourovka Notebook, New Problems, 21st issue (2026), Problem 21.8, with V. G. Bardakov and A. L. Iskra named in the supplied discussion and a definition cross-reference to 17.57. Its background reports affirmative preprints by Pan dated 14 April 2026 and Monticone dated 30 April 2026. The literature triage, checked 2026-08-17, reports Pan's exact all-$k$ theorem and a separate Kourovka project with a formally verified treatment of this problem. It labels the isomorphism question solved with nothing remaining mathematically for the stated assertion, while conventional journal publication of recent work remains pending. The primary registry status is also ``solved''. These are supplied proof-status reports; no paper or formal artifact was independently checked.
\subsection{Short English statement}
Do the horizontal class-transposition groups for moduli $2$ through $k$ together generate the full symmetric action on residue labels modulo their least common multiple? For every $k>3$, the supplied 2026 record reports an affirmative theorem and a separate formally verified treatment.
\subsection{Sources}
\begin{itemize}
\item UnsolvedMath local registry, permanent ID 2517, source code catalogue entry, ``Kourovka Notebook Problem 21.8'', New Problems, 21st issue (2026). The background names V. G. Bardakov and A. L. Iskra and cross-references 17.57 for definitions. Dataset attribution: Ulam AI and the UnsolvedMath Contributors; curation and metadata under CC BY 4.0.
\item Junyao Pan, \emph{A note on the horizontal class transposition group}, arXiv:2604.12553 (2026), preprint recorded as 14 April 2026, \url{https://arxiv.org/abs/2604.12553}.
\item P. Monticone, preprint recorded as 30 April 2026, \url{https://kourovkanotebookorg.wordpress.com/wp-content/uploads/2026/05/21_8-1.pdf}. No title is supplied in the background.
\item Wouter van Doorn, Elias Judin, Pietro Monticone and Daniel Morrison, \emph{On Some Problems from the Kourovka Notebook}, arXiv:2607.17477 (2026), \url{https://arxiv.org/abs/2607.17477}. The supplied assessment cites inclusion of the class-transposition solution among formally verified results.
\item Pietro Monticone et al., Kourovka formalization repository, Problem 21.8 (2026), \url{https://github.com/pitmonticone/Kourovka}. The supplied assessment cites the accompanying indexed formal proof.
\end{itemize}
\end{document}
